\section{Reference architecture}
\label{ch:architecture}

The reference architecture describes how an institutional question becomes an answer that a student can inspect and act on. Its central concern is the connection between evidence and assistance: which source supports a statement, how that source reaches the model, and what happens when the available information is insufficient. It is derived from the chatbot implementation and expressed as responsibilities and contracts that another implementation could retain; its reuse points are proposals grounded in one system.

The description is based on a source review of identified repository revisions. Separate checks matched deployed retrieval and prompt files to the code used for replay. Figure~\ref{fig:architecture} shows the larger application and the smaller evaluated path.

% Native, grayscale-readable source architecture. Required TikZ library: arrows.meta.
% Explicit two/three-line labels avoid paragraph wrapping in boxes.
% Target dimensions: below 17 cm wide and approximately 8.3 cm tall, at 8 pt.
\begin{figure}[htbp]
\centering
\resizebox{\linewidth}{!}{%
\begin{tikzpicture}[
  x=1cm,y=1cm,
  font=\fontsize{8}{9.2}\selectfont,
  box/.style={draw=black!75,line width=.45pt,rounded corners=1.5pt,
    fill=white,align=center,inner xsep=3pt,inner ysep=4pt,
    minimum height=.65cm},
  tested/.style={box,line width=1.25pt,fill=black!6},
  flow/.style={-{Stealth[length=1.6mm,width=1.15mm]},line width=.5pt},
  exchange/.style={arrows={{Stealth[length=1.6mm,width=1.15mm]}-{Stealth[length=1.6mm,width=1.15mm]}},line width=.5pt},
  optional/.style={exchange,dashed},
  note/.style={align=left,inner sep=0pt},
  edgeword/.style={fill=white,inner sep=1.5pt,align=center}
]
  % Preparation: archival originals are distinct from FAQ records.
  \node[note,font=\fontsize{8}{9.2}\selectfont\bfseries,anchor=west]
    at (0,.95) {PREPARATION};
  \node[box,minimum width=1.8cm] (originals) at (.9,0)
    {Originals\\PDF / DOCX};
  \node[box,minimum width=2.05cm] (markdown) at (3.2,0)
    {Markdown\\review required};
  \node[box,minimum width=2.55cm] (ingest) at (6.05,0)
    {Ingestion / worker\\chunk + embed};
  \node[box,minimum width=2.65cm] (vectors) at (9.55,0)
    {PostgreSQL\\pgvector index};
  \node[box,minimum width=3.1cm] (gpu) at (13.75,0)
    {GPU service\\embed + rerank};
  \draw[flow] (originals) -- (markdown);
  \draw[flow] (markdown) -- (ingest);
  \draw[flow] (ingest) -- (vectors);
  \node[note,anchor=south] (faq) at (6.05,.95) {FAQ records};
  \draw[flow] (faq.south) -- (ingest.north);
  \draw[exchange] (ingest.north east) -- (7.5,.85)
    -- node[edgeword,above] {embeddings} (13.75,.85) -- (gpu.north);

  % Serving: reply arrows denote request/stream exchanges, not shared state.
  \node[note,font=\fontsize{8}{9.2}\selectfont\bfseries,anchor=west]
    at (0,-1.05) {SERVING};
  \node[box,minimum width=1.25cm] (web) at (.65,-1.8) {Web UI};
  \node[box,minimum width=1.85cm] (bff) at (2.7,-1.8)
    {Next.js BFF};
  \node[box,minimum width=1.85cm] (discord) at (2.7,-2.9)
    {Discord\\adapter};
  \node[box,minimum width=1.8cm] (api) at (5.5,-2.35)
    {FastAPI\\chat API};
  \node[tested,minimum width=3.1cm,minimum height=1cm] (retrieval)
    at (9,-2.35) {Retrieval + prompt\\selection + context};
  \node[box,minimum width=3.1cm,minimum height=1cm] (agent)
    at (13.75,-2.35) {Provider agent\\hosted / local LLM};
  \draw[exchange] (web) -- (bff);
  \draw[exchange] (bff.east) -- (4.12,-1.8) |- (api.north west);
  \draw[exchange] (discord.east) -- (4.12,-2.9) |- (api.south west);
  \draw[flow] (api) -- (retrieval);
  \draw[flow] (vectors.south) -- (9.55,-1.05) -| (retrieval.north);
  \draw[exchange] (gpu.south west) -- (11.85,-1.05)
    -- (10.85,-1.05) |- (retrieval.north east);
  \draw[flow] (retrieval) -- node[edgeword,above] {prompt} (agent);

  % Experimental transport is an alternative to provider-agent execution.
  \node[box,minimum width=2.7cm] (luna) at (5.5,-4.05)
    {Controlled Luna\\Chat Completions};
  \draw[flow,line width=.8pt] (retrieval.south west)
    -- (7.45,-4.05) -- node[edgeword,above] {replay} (luna.east);

  % Structured records stay outside the FAQ/document embedding path.
  \node[box,minimum width=2.55cm] (localmcp) at (10.1,-4.05)
    {Local MCP\\Python tools};
  \node[box,minimum width=3.05cm] (facultymcp) at (14.35,-4.05)
    {Faculty API / MCP\\Spring AI + JPA};
  \draw[optional] (agent.south west) -- (12.2,-3.3)
    -- (10.1,-3.3) -- (localmcp.north);
  \draw[optional] (agent.south east) -- (15.3,-3.3)
    -- (14.35,-3.3) -- (facultymcp.north);
  \node[box,minimum width=2.2cm] (assets) at (8.2,-5.55)
    {JSON assets\\courses / staff\\rooms / sessions};
  \node[box,minimum width=2.8cm] (adapters) at (11.05,-5.55)
    {Dedicated APIs\\EduPage timetables\\faculty thesis lists};
  \node[box,minimum width=3.05cm] (facultydb) at (14.35,-5.55)
    {Faculty PostgreSQL\\operational records};
  \draw[flow] (assets.north) -- (8.2,-4.75) -| (localmcp.south);
  \draw[flow] (adapters.north) -- (11.05,-4.75) -| (localmcp.south);
  \draw[flow] (facultydb) -- (facultymcp);

  \node[note,text width=16.35cm,anchor=north west] at (0,-6.55)
    {\textbf{Bold border:} production retrieval/prompt code tested. \enspace
     \textbf{Dashed arrows:} optional MCP, not evaluated.\\
     Controlled Luna replaces the provider agent. Storage, stream resumption and monitoring are omitted.};
\end{tikzpicture}}
\caption{Chatbot architecture and the components tested.
Structured-data lookup is separate from document/FAQ retrieval. Reading static feeds does not establish how those feeds are published.
The diagram shows features found in the inspected code; not all were verified in deployment.}
\label{fig:architecture}
\end{figure}


\subsection{Responsibilities and architectural contracts}

The system has six connected responsibilities: preparing sources, retaining knowledge, selecting evidence, constructing and generating answers, delivering conversations, and recording state and observations. Structured tools provide an additional evidence path for records that are better queried through an API than represented as document passages. These responsibilities need not map one-to-one to machines or repositories. Separating them makes individual decisions inspectable: a poor answer can originate in an incorrect conversion, missed evidence, an unsupported inference, or information lost during delivery.

Table~\ref{tab:architecture-contracts} describes the principal boundaries. A contract here means the information and obligations crossing a boundary, not a claim that every obligation is enforced automatically. For example, an ingestion endpoint can require a document name while the accuracy of the contents still depends on editorial review, so a successful database write is not validated knowledge.

\begin{table}[htbp]
\centering
\caption{Principal contracts in the reference architecture.}
\label{tab:architecture-contracts}
\small
\begin{tabularx}{\linewidth}{@{}p{0.20\linewidth}XX@{}}
\toprule
Boundary & Information transferred & Architectural obligation \\
\midrule
Source to ingestion & Stable name, title, text and provenance & Preserve a connection to the original and distinguish conversion from authority. \\
Knowledge to retrieval & FAQ records or ordered document chunks and current vectors & Keep source type and parent identity available during selection. \\
Retrieval to answering & Selected text, source identifiers, links and context status & Expose the actual evidence supplied to generation. \\
Tools to answering & Structured records, ambiguity or lookup failure & Distinguish a returned record from permission to take an institutional action. \\
Answering to channel & Response identity and streamed events & Preserve answer meaning and communicate completion or failure. \\
Request to observation & Request/response identifiers and stage outcomes & Separate operational traces, mutable state and evaluation evidence. \\
\bottomrule
\end{tabularx}
\end{table}

A successful component does not guarantee a good answer. Retrieval can return a regulation that has been superseded; a provider can finish generating text without covering all conditions of a procedure; a channel can display the answer without a separate source event. End-to-end effectiveness therefore requires judging whether the final assistance is supported and sufficiently complete, or provides an appropriate next step.

\subsection{Source provenance and the knowledge lifecycle}

FAQs, regulations, course and timetable records, and historical correspondence play the different roles described in Chapter~\ref{ch:data}, and the architecture keeps these roles explicit instead of treating every available text as interchangeable knowledge.

The document lifecycle begins with an original file and a derived textual representation. Retaining both lets an editor compare a conversion with its source and correct it later. Ingestion assigns the document a stable identity, stores its title and provenance, and divides its text into chunks that remain connected to their parent and position. This matters when a retrieved sentence depends on a heading or a preceding condition, and it lets an answer refer to an identifiable document rather than an anonymous fragment.

Provenance has several independent dimensions: original-file identity, conversion quality, source authority, and currentness. A hash can detect changed bytes but establishes none of the latter three, and an authority URL supports inspection without proving applicability to the student's situation. Currentness metadata must therefore be supplied and interpreted, not inferred from recent ingestion. The inspected application can expose status, date and verification fields when they exist; reuse at another institution would require explicit ownership of those decisions.

Vector freshness is a separate technical concern. If indexed text changes while its vector does not, retrieval can reflect an earlier version. The implementation addresses this through invalidation and revision-aware embedding updates, which keep stored content and its derived representation consistent without determining whether the rule is correct. Optional object storage serves archival access to originals, while PostgreSQL stores retrieval records.

\subsection{Retrieval as an inspectable evidence-selection stage}

Retrieval operates over two corpora: short FAQ records and passages from longer documents. Dense search finds candidates in each corpus, optional lexical expansion contributes FAQ matches under the application's conditions, and a cross-encoder reranks the combined candidates. The selected candidates retain their source type, identity and score, so selection policy can be varied while inspecting exactly which evidence enters the answer prompt.

The shared candidate pool creates a substantive design choice. A FAQ may restate the user's question, while a regulation may contain the relevant condition in formal language, so equal reranker scores need not have equal implications for the two source types. A common threshold is simple, but lowering only the document threshold can admit weakly related passages. The architecture therefore exposes thresholds as variation points.

Selection also includes allocation and expansion. The inspected path gives FAQ candidates priority within the final candidate budget and can add neighboring passages around selected document chunks. A limit on selected candidates is therefore not a fixed limit on characters presented to the model; one long document section can change the context more than one short FAQ. Comparisons must preserve or report these decisions, or an apparent threshold effect may include changes in ordering or context construction.

Evidence admission is distinct from source presentation. The inspected default context threshold is 0.10, while the source-display threshold is 0.30, so a passage may influence generation without producing a visible source card. Retrieval and a useful-link catalog are also assembled separately: empty retrieved text may still be accompanied by links that help the model suggest a relevant destination. Recording text and links separately allows useful routing to be distinguished from a supported substantive answer.

\subsection{Answer construction and structured lookup}

Answer construction combines the question, selected evidence, useful links, date, available history and interface instructions. Provider adapters translate this input into the appropriate generation route and emit a common response stream. The generated answer may summarize or explain the retrieved material, but remains an interpretation that must be assessed for support and completeness.

Structured tools extend this process. The inspected agent can discover configured MCP tools and call them during generation, for example to look up named courses, rooms, schedules or catalog records. This access pattern differs from semantic search over a regulation, and its failures differ: the corpus may lack a procedural rule, whereas a structured service may be unavailable or return an ambiguous entity match.

Read-only lookup does not resolve an individual case. A thesis listing or professor record may inform guidance, but does not appoint a committee, approve a request or register a student. Service authentication allows one system to contact another; it does not establish the student's authorization to inspect protected records. The architecture therefore separates explaining a procedure, locating public or service-accessible information, and exercising institutional authority.

Clarification and referral are valid outcomes when the evidence cannot support a direct answer. A general procedure may depend on the academic period or the type of request, and asking for that distinction may be more useful than presenting one branch as universal. A referral can identify the appropriate office while leaving the decision to that authority.

\subsection{Channels, state and failure boundaries}

Web and Discord adapters send requests to the same backend but maintain different conversation and delivery behavior. The web path has persisted conversations and resumable streams. Discord retains bounded process-local history and edits or splits messages to fit its presentation limits. Shared generation reduces duplicated backend logic, but histories, source displays and interruption behavior can differ between the two.

State is separated by purpose. Knowledge records represent evidence; conversations represent mutable user interaction; lookup caches represent previously fetched data; telemetry represents observations; frozen research artifacts represent the inputs and outputs of a particular study. Merging them would obscure retention, authority and reproducibility. Response identifiers can connect events and feedback without making either a substitute for content assessment.

Failure boundaries make the system easier to interpret. A request may fail admission before retrieval. Retrieval may fail to obtain reranker scores. A tool refresh may fail while an older inventory remains available. Generation may produce partial text before an error, and a client may lose continuity although the backend completed. These outcomes require different evidence and remedies; in particular, serving old cached information preserves availability but cannot be reported as a successful refresh.

\subsection{Variation points and empirical coverage}

An institution adopting this architecture would replace its corpus, source-review process, terminology, structured adapters and identity policy, while retaining the separation of source preparation, evidence selection, answer construction and channel delivery. Embedding models, rerankers, candidate budgets, selection thresholds, generation providers and presentation rules are additional variation points whose value depends on local questions and evidence.

The experiment covers a deliberately smaller boundary. It reuses production retrieval and prompt construction with the original queries, then sends captured prompts to controlled Luna generation. It excludes query rewriting, conversation history, live provider agents, MCP and built-in tool execution, request-admission checks, conversation persistence, and actual web or Discord rendering. The input schema still affects eligibility, as documented in the methodology. Evidence selection thus changes while other captured inputs remain controlled, which makes the retrieval comparison interpretable; the remaining responsibilities require their own validation.
